\documentclass[main.tex]{subfiles}
\begin{document}
\label{conclusion}

In conclusion, this study found that during the COVID-19 pandemic restrictions, the mobility of individuals collapsed into small geographic areas. The geometric distribution of these trips remains constant as they shrink, highlighting two main locations that account for approximately 70$\%$ of all visited positions. These locations serve as the starting points for discovering new places. Notably, discovery trips tend to be shorter, as evidenced by the change in the exponent of the tail of the probability distribution of trip lengths. Despite this, the relative frequency of longer trips remained unchanged since they represent a small fraction of all the trips. This evidence, combined with the fact that people's exploration was reduced, shows that a relatively big portion of exploration is covered by short trips. Testing different entropies in general during COVID, the information is lower for the population; this fact contributes to the lack of number of points in trajectory, but also to the lack of structure.

Looking at the last figure in \ref{entropic_measures_california}, it seems to be the case that the effect on the lowering of the entropy is due more to the lack of points rather the lack of structure. This suggests that the higher content of structure and therefore in information is concentrated in smaller trips that therefore are more varied. In the end, conditioning the study of the radius of gyration to the government party across different counties and to the level of rurality, one can see that if on the one hand the government party is not an indicator of how much people will move (lockdowns are essentially the same), rurality shows that people feel more free to move if they are in rural areas. The post-lockdown period shows that people coming from democratic areas tend to decrease their willingness to move compared to republicans.

Summarising, in this work, we explored the changes in individual mobility habits under COVID-19 non-pharmaceutical intervention. Our analysis provides a set of different measures to appreciate the changes due to lockdowns that, according to our observation, strongly influenced individual mobility patterns not only by restricting the number of visits, the diversity of the activities done at the destination, or their duration~\cite{lucchini}, but also strongly modifying the spatial characteristics of the individual mobility networks and the natural attitude of people of exploring new locations.


\end{document}


